\documentclass[../NDCNNAI_main.tex]{subfiles}
\graphicspath{{\subfix{../fig/}}}

\begin{document}
\subsection{Определение и декомпозиция ошибки DeepONet}

Пусть \( X = C(D) \) --- пространство непрерывных функций на компакте \( D \subset \mathbb{R}^d \), а \( G \colon X \to L^2(U) \) --- целевой оператор, который мы хотим аппроксимировать. Пусть \( \mu \) --- вероятностная мера на \( X \) (распределение входных функций \( u \)).

\begin{definition}[Среднеквадратичная операторная ошибка DeepONet]
Для DeepONet \( N \colon X \to L^2(U) \) определим ошибку:
\[
\mathcal{E}(N) \coloneq \|G - N\|_{L^2(\mu; L^2(U))} 
= \left( \int_X \|G(u) - N(u)\|_{L^2(U)}^2 d\mu(u) \right)^{1/2}
\]
\end{definition}

\begin{remark}[1]
    \( u \sim \mu \) --- случайная входная функция
\end{remark}
\begin{remark}[2]
    \( G(u) \) и \( N(u) \) --- случайные выходные функции на \( U \)
\end{remark}
\begin{remark}[3]
    \( \mathcal{E}(N)^2 \) --- математическое ожидание квадрата \( L^2(U) \)-ошибки между \( G(u) \) и \( N(u) \)
\end{remark}
\begin{remark}[4]
    На практике эмпирическая оценка среднеквадратичной ошибки является оценкой Монте-Карло для \( \mathcal{E}(N)^2 \) по конечной выборке \( \{u^{(i)}\}_{i=1}^n \)
\end{remark}

Стандартная архитектура DeepONet: \( N = R \circ A \circ E \), где
\begin{itemize}
    \item \( E \colon X \to \mathbb{R}^m \) --- кодировщик (кодирует входную функцию в вектор),
    \item \( A \colon \mathbb{R}^m \to \mathbb{R}^p \) --- аппроксиматор (нейронная сеть, отображает вектор в коэффициенты),
    \item \( R \colon \mathbb{R}^p \to L^2(U) \) --- реконструктор (строит выходную функцию по коэффициентам).
\end{itemize}

Для такой архитектуры ошибку можно разложить на три компоненты:

\begin{equation}
\mathcal{E}(N) \leq 
\underbrace{\mathcal{E}_{\text{enc}}}_{\text{ошибка\\кодирования}} + 
\underbrace{\mathcal{E}_{\text{app}}}_{\text{ошибка\\аппроксимации}} + 
\underbrace{\mathcal{E}_{\text{rec}}}_{\text{ошибка\\реконструкции}}
\end{equation}
Если подробнее:
\begin{itemize}
    \item \textbf{\( \mathcal{E}_{\text{enc}} \)} --- потеря информации при переходе от функции \( u \) к конечномерному вектору \( E(u) \).
    \item \textbf{\( \mathcal{E}_{\text{app}} \)} --- ошибка нейронной сети \( A \), аппроксимирующей идеальное отображение \( u \mapsto (\beta_1(u), \dots, \beta_p(u)) \).
    \item \textbf{\( \mathcal{E}_{\text{rec}} \)} --- ошибка, связанная с использованием конечного базиса \( \{\tau_k\}_{k=1}^p \) для реконструкции и игнорирования высокочастотных мод.
\end{itemize}

\begin{example}
Приведём примеры оценок погрешностей кодирования и реконструкции.
\begin{itemize}
    \item \textbf{Случайные сенсоры} Для случайных полей с аналитической ковариацией (характеризуется величиной \( \ell \)) и равномерных сенсоров с декодером на основе Фурье:
    \[
    \mathcal{E}_{\text{enc}} \lesssim \exp\big(-c\ell m^{1/d}\big)
    \]
    \item \textbf{Спектральные реконструкторы} Если выбрать \( \{\tau_k\}_{k=1}^p \) как собственные функции ковариации \( G\#\mu \) или как базис Лежандра/Фурье (при соболевской гладкости \( k \)), то:
    \[
    \mathcal{E}_{\text{rec}} \lesssim p^{-k/d}
    \]
\end{itemize}
В сочетании с голоморфной зависимостью параметров и стандартными оценками для аппроксиматора $A$ даёт явные двусторонние оценки.    
\end{example}


\subsection{Измеримая УАТ для DeepONet}

\begin{theorem}[Измеримая УАТ для DeepONet]
Пусть \( \mu \) --- вероятностная мера на $C(D)$ и \( G \colon C(D) \to L^2(U) \) --- измеримый оператор, причём \( G \in L^2(\mu; L^2(U)) \). Тогда для любого \( \varepsilon > 0 \) существует DeepONet \( N = R \circ A \circ E \) такой, что
\[
\mathcal{E}(N) = \|G - N\|_{L^2(\mu; L^2(U))} < \varepsilon
\]
\end{theorem}

\begin{proof}
Опишем общую схему доказательства:
\begin{enumerate}
    \item По теореме Лузина для любого \( \delta > 0 \) существует компакт \( K \subset C(D) \) с \( \mu(K) > 1 - \delta \) такой, что \( G|_K \) непрерывен.
    \item Используя теорему Чен--Чен (аналог УАТ для непрерывных операторов), строим DeepONet \( N_K \), равномерно приближающий \( G \) на \( K \):
    \[
    \sup_{u \in K} \|G(u) - N_K(u)\|_{L^2(U)} < \varepsilon/2.
    \]
    \item Полагаем \( N(u) = N_K(u) \) для \( u \in K \) и, например, \( N(u) = 0 \) для \( u \notin K \). Тогда $N$ является борелевским измеримым оператором
    \[
    \mathcal{E}(N)^2 = \int_K \|G - N\|^2 d\mu + \int_{K^c} \|G\|^2 d\mu \leq (\varepsilon/2)^2 + \int_{K^c} \|G\|^2 d\mu.
    \]
    Поскольку \( G \in L^2(\mu) \), второе слагаемое можно сделать меньше \( (\varepsilon/2)^2 \) выбором достаточно малого \( \mu(K^\complement) \).
\end{enumerate}
\end{proof}

\subsection{<<Upper bound program>>: получения верхних оценок}
Теорема об измеримом УАТ обосновывает следующий конструктивный подход к оценке ошибки DeepONet:
\begin{enumerate}
    \item Оценка \( \mathcal{E}_{\text{enc}} \): проектирование эффективных кодировщиков/декодеров \( (E, D) \), минимизирующих потерю информации.
    \item Оценка \( \mathcal{E}_{\text{rec}} \): выбор базиса реконструкции \( R \), оптимального относительно спектральных свойств распределения \( G\#\mu \).
    \item Оценка \( \mathcal{E}_{\text{app}} \): применение стандартных теорем аппроксимации нейронных сетей в конечномерных пространствах.
\end{enumerate}

\subsection{Применение к динамическим системам: от временных рядов к операторам}

Рассмотрим нелинейную динамическую систему.

\textbf{Дискретное время.}
\[
x_{t+1} = F(x_t, u_t), \quad t = 0, 1, 2, \dots,
\]
где \( x_t \in \mathbb{R}^{d_x} \) --- состояние, \( u_t \in \mathbb{R}^{d_u} \) --- управление.

\textbf{Непрерывное время (ОДУ).}
\[
\dot{x}(t) = f(x(t), u(t), t), \quad x(a) = x_0, \quad t \in [a, b].
\]
Вместо пошагового предсказания можно рассматривать оператор решения:
\[
G \colon u(\cdot) \longmapsto x(\cdot),
\]
который ставит в соответствие входному сигналу \( u(t) \) соответствующую траекторию \( x(t) \).
Для автономных систем также рассматривается оператор временного сдвига
\[
S_h \colon x(\cdot) \longmapsto x(\cdot + h),
\]
который сдвигает траекторию по $h$.

\begin{example}[Оператор решения ОДУ]
Рассмотрим систему на \([a, b]\):
\[
\begin{cases} 
\frac{ds}{dx} = g(s(x), u(x), x), \\
s(a) = s_0 \in \mathbb{R}^K,
\end{cases}
\]
где \( u \in V \subset C([a, b]) \) --- входной сигнал. Оператор решения:
\[
(Gu)(x) = s_0 + \int_a^x g((Gu)(t), u(t), t) \, dt.
\]

\emph{Задача идентификации}: по наблюдениям пар \( (u, s) \) построить DeepONet \( N \), аппроксимирующий \( G \).

\textbf{Дискретизация входа: сенсоры и интерполяция.}
Выбираем \( m+1 \) равномерных сенсоров на \([a, b]\):
\[
x_j = a + j \frac{(b - a)}{m}, \quad j = 0, 1, \dots, m.
\]
Строим кусочно-линейный интерполяцию \( u_m \). При компактности \( V \) и условии Липшица на \( g \) выполняется оценка в равномерной норме:
\[
\max_{x \in [a, b]} |u(x) - u_m(x)| \leq \kappa(m, V) \to 0 \quad \text{при } m \to \infty.
\]

\textbf{Оценка необходимого числа сенсоров.}
Если \( g \) удовлетворяет условию Липшица по \( (s, u) \), то по неравенству Гронвала имеет место оценка в норме $L^2$ или $L^{\infty}$ по времени:
\[
\|(Gu)(x) - (Gu_m)(x)\|_2 \leq c(b - a) \kappa(m, V) e^{c(b - a)}, \quad x \in [a, b].
\]

Таким образом, выбрав \( m \) так, чтобы 
\[
c(b - a) \kappa(m, V) e^{c(b - a)} < \varepsilon,
\]
мы гарантируем \( \varepsilon \)-близость решений в норме $L^2$ или $L^{\infty}$ на $[a, b]$. После этого достаточно аппроксимировать конечномерное отображение
\[
(u(x_0), \dots, u(x_m)) \mapsto (Gu)(d)
\]
обычной нейронной сетью.
\end{example}

В конечном счёте можно предложить следующею модель применения DeepONet к нелинейным динамическим системам, применительно к \textbf{branch/кодировщик (входная часть):}
\begin{itemize}
    \item Вход: дискретизированное управление или состояние на коротком временном окне, например $(u(t_1), \dots, u(t_m))$ или $(x(t_1), \dots, x(t_m))$.
    \item Кодировщик $E$ отображает это в $\mathbb{R}^m$ (частно просто тождественное отображение).
    \item Аппроксиматор $A\colon \mathbb{R}^m \to \mathbb{R}^p$ является глубокой MLP, выдающей коэффициенты $\alpha = (\alpha_1, \dots, \alpha_p)$, зависящие лишь от окна.
\end{itemize}
Со стороны выхода \textbf{trunk/реконструктор:}
\begin{itemize}
    \item Trunk вход: время запроса $t \in [a, b]$ (или пространственная переменная $y$).
    \item Trunk сеть базисные функции $\tau_k(t)$ и выходы
    $$(N(u))(t) = \tau_0(t) + \sum_{k = 1}^{p}\alpha_k\tau_k(t).$$
    \item Вычисление $t$ на сетке даёт приближение всей траектории.
\end{itemize}
Для обучения можно брать пары $(u^{(i)}, x^{(i)}(\cdot))$ или перемещать окно по длинным траекториям, а в качестве функции потерь --- эмпирическое приближение $\mathcal{E}(N)$
$$
\frac{1}{N_{\mathrm{data}}}\sum_{i}\int_{a}^{b}\|x^{(i)}(t) - (Nu^{(i)})(t)\|_2^2\, dt.
$$
\subsection{ISO/BIBO-ориентированная регуляризация}
Когда используются DeepONet для динамических систем, часто важна \emph{устойчивость} (устойчивость состояния по выходу, сохранение ограниченности выхода при наличии таковой у входа).
Добиться этого можно, накладывая \textbf{штрафы на основе констант Липшица или Якобиана}
\begin{itemize}
    \item Аппроксиматор $A$ определяет чувствительность коэффициентов $\alpha$ к закодированному входу $E(u)$.
    \item Мы можем штрафовать во время обучения за слишком большую операторную норму $A(E(u))$, чтобы малость константы Липшица $A \circ E$ и изменений в окне приводила к незначительны изменениям траекторий.
    \item Это аналогично ограничениям типа Ляпунова/ISS в классической модели и может помочь избежать взрыва предсказаний при экстраполяции.
\end{itemize}

Таким образом в отличии от классического подхода, где нужно отдельно обучать модель для каждой точки, DeepONet делает эту конструкцию явной:
\begin{itemize}
    \item Кодировщик \(E\): вычисляет значения входной функции \(u\) в точках сенсоров.
    \item Аппроксиматор \(A\): обучается предсказывать значения решения для множества запрашиваемых точек на основе данных сенсоров.
    \item Реконструктор \(R\): собирает предсказания для всех точек запроса в полную траекторию на отрезке \([a, b]\).
\end{itemize}

Для углубления в этот самый свежий раздел курса рекомендуется обратится к первоисточникам \cite{Lu2021, Telgarsky2016, Hanin2019, Lanthaler2021}.
\end{document}
